% Clase 22 --- La impedancia en el plano complejo (§3.3).
% Ésta es la figura que hace que todo el método complejo se vuelva geométrico:
% el módulo de Z da la amplitud y su argumento el desfasaje. Los dos paneles
% son el mismo triángulo con el signo de (omega L - 1/omega C) invertido.
\begin{center}
\begin{tikzpicture}[scale=1]

  % =============== (a) inductivo ===============
  \begin{scope}
    \draw[figaxis] (-0.6,0) -- (2.9,0);
    \draw[figaxis] (0,-1.9) -- (0,1.9);
    \node[figlblaux, anchor=north east] at (2.9,-0.05) {Re};
    \node[figlblaux, anchor=south east] at (-0.08,1.9) {Im};

    \draw[figvec, line width=1.2pt] (0,0) -- (2.0,0);
    \node[figlbl, figblue, anchor=north] at (1.0,-0.1) {$R$};
    \draw[figvec, line width=1.2pt] (2.0,0) -- (2.0,1.35);
    \node[figlbl, figblue, anchor=west] at (2.08,0.7)
      {$\omega L - \dfrac{1}{\omega C}$};
    \draw[figvecres, line width=1.4pt] (0,0) -- (2.0,1.35);
    \node[figlbl, figred, anchor=south east] at (1.15,0.78) {$\hat Z$};
    \draw[figaux] (0.75,0) arc (0:34:0.75);
    \node[figlblaux, anchor=west] at (20:0.95) {$\varphi>0$};

    \node[figlbl, anchor=north, align=center] at (1.2,-2.35)
      {(a) $\omega L > \dfrac{1}{\omega C}$: inductivo,\\[-0.1em]
       la corriente \textbf{atrasa}};
  \end{scope}

  % =============== (b) capacitivo ===============
  \begin{scope}[xshift=7.4cm]
    \draw[figaxis] (-0.6,0) -- (2.9,0);
    \draw[figaxis] (0,-1.9) -- (0,1.9);
    \node[figlblaux, anchor=north east] at (2.9,-0.05) {Re};
    \node[figlblaux, anchor=south east] at (-0.08,1.9) {Im};

    \draw[figvec, line width=1.2pt] (0,0) -- (2.0,0);
    \node[figlbl, figblue, anchor=south] at (1.0,0.08) {$R$};
    \draw[figvec, line width=1.2pt] (2.0,0) -- (2.0,-1.35);
    \node[figlbl, figblue, anchor=north west] at (2.08,-1.05)
      {$\omega L - \dfrac{1}{\omega C}$};
    \draw[figvecres, line width=1.4pt] (0,0) -- (2.0,-1.35);
    \node[figlbl, figred, anchor=north east] at (1.15,-0.78) {$\hat Z$};
    \draw[figaux] (0.75,0) arc (0:-34:0.75);
    \node[figlblaux, anchor=west] at (-20:0.95) {$\varphi<0$};

    \node[figlbl, anchor=north, align=center] at (1.2,-2.35)
      {(b) $\omega L < \dfrac{1}{\omega C}$: capacitivo,\\[-0.1em]
       la corriente \textbf{adelanta}};
  \end{scope}

\end{tikzpicture}\\[0.5em]
{\small\itshape El truco del método complejo es que, al proponer
$\hat I = \hat I_M e^{j\omega t}$, cada derivada se convierte en multiplicar por
$j\omega$ y la ecuación diferencial se vuelve una división de números
complejos: $\hat I_M = \varepsilon_0/\hat Z$, con
$\hat Z = R + j(\omega L - 1/\omega C)$. Toda la información del circuito queda
en ese único número, y la figura lo dice mejor que la fórmula: la parte real es
la resistencia ---lo que disipa--- y la parte imaginaria es la
\textbf{reactancia}, la competencia entre bobina y condensador, que no disipa
pero desfasa. El módulo de la hipotenusa fija la amplitud,
$I_M = \varepsilon_0/|\hat Z|$, y el ángulo fija el desfasaje entre corriente y
tensión. Según cuál de los dos gane, el circuito se comporta como inductivo o
como capacitivo. Y hay un caso especial que se lee de un vistazo: cuando los dos
se cancelan, el triángulo se aplasta sobre el eje real, $\hat Z = R$, y estamos
en resonancia.}
\end{center}
